\begin{afterword}
\summaryhead

The post defines the halo effect as the affect heuristic at work in judgments of people: one good trait raises the judgment of others. Its evidence is a long passage from Robert Cialdini's \textsc{Influence}. Attractive people, Cialdini reports, are credited with talent, kindness, honesty and intelligence. Attractive candidates won far more votes in a Canadian election, while voters denied being swayed. Well-groomed applicants were favored in hiring, attractive workers earn more, attractive defendants were twice as likely to avoid jail, and attractive people get more help and persuade more.

The author calls the effect of looks on judgments of honesty a clear bias, since the two have no legitimate connection, while granting that honesty and intelligence may really go together. The studies should make readers suspect a similar halo among kindness, intelligence and other virtues, and distrust a world that sorts too neatly into devils and angels. The post ends by advising a little more skepticism toward attractive political candidates.

\responsehead

The post's main claim, that attractiveness colors judgments of unrelated traits, is well supported, and its closing advice follows from real evidence. The problems are in how the evidence is presented.

All of it comes through one passage from a popular book, and the passage is weakest where the post relies on it most. Cialdini lists kindness and honesty among the traits attractive people are credited with, citing a meta-analysis. That review found the effect near zero for integrity and concern for others, and moderate on average. The post's next paragraph takes honesty as its clearest case of bias. The voting figure comes, by Cialdini's own notes, from an unpublished manuscript; published work since then finds the same direction at a smaller size. The jail figure is a raw comparison, in a study whose attractive defendants also faced less serious charges, and whose author called the relationship weak.

The post's own reasoning is sound. It sees that some traits really do go together, so that not every correlation in judgments is a bias. Its extension, that the attractiveness studies ``should make us suspicious that there may be a similar halo effect for kindness, or intelligence,'' has stronger support than the post gives it. Thorndike named the halo effect in 1920 for ratings of intelligence, industry, reliability and character that correlated more than reality could explain, and he drew the same distinction the post draws between real correlation and inflated correlation. What the post offers as a suspicion was the field's starting point.

In short: a sound point about looks and judgment, argued from a popular summary that in places claims more than its sources, and a sound extension that the original research on the halo effect supports.
\end{afterword}
